\subsection*{一步一步建立直观}

\paragraph{1. 用一个能亲手改变的对象开始}

先把两根箭头当作地图上的两次移动，拖动它们的端点，描述合位移如何变化。此处不需要预先知道向量空间的定义；目标是先形成一个明确问题：哪些变化来自移动的次序，哪些来自方向和长度？这会成为后面比较不同表示方法的练习。



\paragraph{2. 每次只改变一件事，再解释结果}

保持两根箭头不动，只改变其中一根的系数。先写下你预计的方向和长度，再查看结果。学习数学时，操作的价值在于让预测暴露出来，而不是替代思考。若结果不符，应追问是计算错误、坐标误解，还是原先使用的规则缺少条件。



\paragraph{3. 把特殊情形当作检查工具}

让两根向量共线，观察原先能覆盖平面的组合怎样退化到一条线。零、对称、相等和极端情形经常能揭示一般公式隐藏的条件。本条目是学习路线说明；右图只是这一学习方法的示例，不把课程公告改写成一条不存在的原课数学定理。
